\section{Categorical framework}\label{sec:categorical-framework}
%%% ====================================================================

The preceding sections established that $G\colon\Ztwenty\to\mathbb{Z}_2\times\mathbb{Z}_2$
determines both the Euler product of $\zeta_{\mathbb{Q}(\sqrt{5})}(s)$ (upper branch of the
diagram in \cref{sec:objective}) and the spectral structure
$\mu_3=1/2$, $S_3$, $\omega^k/2$ (lower branch).
We now make the categorical content of this determination explicit:
the Frobenius system is a functor (\cref{ssec:frobenius-functor}),
the Euler product is a colimit (\cref{ssec:euler-colimit}),
the Dedekind factorisation is a natural
isomorphism (\cref{ssec:dedekind-natural}),
and the two branches form a span with common
source (\cref{ssec:span}).
Throughout, we indicate what is invocable from the Primitive
Structure~\cite{V11} and what is developed here.

% ------------------------------------------------------------------
\subsection{Preliminaries}\label{ssec:cat-prelim}
% ------------------------------------------------------------------

\begin{definition}[Category of Dirichlet series]\label{def:DirSer}
The category $\mathbf{DirSer}$ has as objects the Dirichlet series
$\sum a_n n^{-s}$ absolutely convergent in some half-plane
$\mathrm{Re}(s)>\sigma_a$, with the Dirichlet convolution
$(a*b)_n = \sum_{d|n} a_d\,b_{n/d}$ as monoidal product and
$\zeta_{\mathrm{triv}}(s)=1$ as unit.
For the partial Euler products below, the relevant morphisms are the
inclusion maps $E(S)\hookrightarrow E(S')$ for $S\subseteq S'$,
which add factors without altering existing coefficients.
\end{definition}

\begin{definition}[Spectral category]\label{def:SpecOmega}
$\mathrm{Spec}(\OmH)$ is the category whose objects are spectral
triples $(\sigma,\mu,\{\lambda_k\})$ satisfying $\sigma+\mu=2$,
$\sigma\mu=3/4$, $\lambda_k = (\mu/1)\omega^k$.
By spectral uniqueness (\cref{cor:S3-from-euler}), the only
solution with $\mu<1$ is $(\sigma,\mu)=(3/2,1/2)$; hence
$\mathrm{Spec}(\OmH)$ is the terminal category (one object, one
morphism): the condition $\mu<1$ that excludes the degenerate solution
$(\sigma,\mu)=(1,1)$ is the No-Diagonal theorem
(Theorem~3 of~\cite{V11}), which establishes
$\|\Omega\|=1/2<1$ (\cref{prop:worldline-entropy}).
The functor $F_{\mathrm{spec}}\colon
\mathbf{B}\Ztwenty\to\mathrm{Spec}(\OmH)$ collapses all eight
morphisms of $\mathbf{B}\Ztwenty$ to this unique spectral point.
\end{definition}

% ------------------------------------------------------------------
\subsection{The Frobenius system as a functor}\label{ssec:frobenius-functor}
% ------------------------------------------------------------------

\begin{definition}[Frobenius functor]\label{def:frobenius-functor}
The \emph{Frobenius functor} of $\Rpcf$ is the monoid homomorphism
\begin{equation}\label{eq:Psi-functor}
  \Psi\colon (\mathbb{N}_{>0},\times)
  \longrightarrow \mathrm{End}_{\mathrm{Ring}}(\Rpcf),
  \qquad
  \Psi(n)(\golden) = \golden^n = F_n\,\golden + F_{n-1}.
\end{equation}
\end{definition}

\begin{proposition}[Functoriality of $\Psi$]\label{prop:Psi-functorial}
$\Psi$ preserves composition and identity:
$\Psi(mn) = \Psi(m)\circ\Psi(n)$ and $\Psi(1)=\mathrm{id}_{\Rpcf}$.
The restriction to primes generates $\Psi$ through
$\Psi(p_1^{a_1}\cdots p_k^{a_k}) = \psi_{p_1}^{a_1}\circ\cdots\circ\psi_{p_k}^{a_k}$.
\end{proposition}

\begin{proof}
$\Psi(mn)(\golden) = \golden^{mn} = (\golden^m)^n
= \Psi(n)(\Psi(m)(\golden))$.
Since $\golden$ generates $\Rpcf$ over $\mathbb{Z}[\tfrac{1}{2}]$,
the endomorphism is determined by its action on $\golden$.
\end{proof}

\begin{remark}[$\Lambda$-ring structure]\label{rmk:lambda-ring}
In the language of Borger~\cite{Borger}, $\Psi$ restricted
to primes equips $\Rpcf$ with a $\Lambda$-ring structure:
$\psi_p(a)\equiv a^p\pmod{p}$ and $\psi_p\circ\psi_q = \psi_{pq}$,
constituting $\mathbb{F}_1$-descent data.
The functor $\Psi$ is the algebraic backbone of the chain
$\golden^2=\golden+1 \to \psi_p \to \chi_5 \to G \to \prod f_p$.
\end{remark}

\begin{lemma}[From Frobenius to Euler factors]\label{lem:psi-to-euler}
The partial Euler functor $E$ (\cref{def:partial-euler})
is determined by $\Psi$ through the composition
$\Psi(p)(\golden) = \golden^p \to F_p\bmod p \to \chi_5(p) \to f_p(s)$:
concretely, $\chi_5(p)\equiv F_p\pmod{p}$
(\cref{prop:fib-split}) and $f_p$ is determined by
$\chi_5(p)$ (\cref{prop:local-factors}).
\end{lemma}

% ------------------------------------------------------------------
\subsection{The Euler product as a categorical colimit}%
\label{ssec:euler-colimit}
% ------------------------------------------------------------------

\begin{definition}[Partial Euler products]\label{def:partial-euler}
Let $\mathcal{P}$ denote the category whose objects are finite sets
$S$ of primes and whose morphisms are inclusions $S\hookrightarrow S'$.
The \emph{partial Euler functor} is
\begin{equation}\label{eq:partial-euler}
  E\colon \mathcal{P} \longrightarrow \mathbf{DirSer},
  \qquad
  E(S)(s) = \prod_{p\in S} f_p(s),
\end{equation}
where $f_p(s)$ is the local factor determined by $\chi_5(p)$
(\cref{prop:local-factors}).
Each $f_p$ depends only on $\chi_5(p) = \pi_2(G(p\bmod 20))$
(\cref{thm:bridge3}).
\end{definition}

\begin{proposition}[Euler product as colimit]\label{prop:euler-colimit}
For $\mathrm{Re}(s)>1$, the Dedekind zeta function is the colimit
of the partial Euler functor:
\begin{equation}\label{eq:euler-colimit}
  \zeta_{\mathbb{Q}(\sqrt{5})}(s) = \varinjlim_{S\in\mathcal{P}} E(S)(s).
\end{equation}
The colimit exists and is unique in $\mathbf{DirSer}$: each
coefficient $a_n$ depends only on the primes dividing $n$ and
stabilises once $S$ contains those primes.
\end{proposition}

\begin{proof}
For each $n\in\mathbb{N}_{>0}$, the $n$-th coefficient of $E(S)(s)$
equals $a_n$ whenever $S$ contains all prime factors of $n$.
The system $\{E(S)\}_{S\in\mathcal{P}}$ is coherent:
inclusions $S\subseteq S'$ add factors without altering existing
coefficients.  The colimit is the unique Dirichlet series
whose $n$-th coefficient agrees with $a_n$ for all $n$.
Convergence for $\mathrm{Re}(s)>1$ ensures existence as an
analytic object.
\end{proof}

\begin{remark}[Monoidal structure]\label{rmk:monoidal}
The colimit is monoidal: each $E(S) = \bigotimes_{p\in S} f_p$
in the Dirichlet convolution, and the colimit is the infinite
$\bigotimes$.  This connects to the ObjectNorm of the Primitive
Structure~\cite{V11}: the multiplicative norm
$\norm{X\otimes Y} = \norm{X}\cdot\norm{Y}$ is the same
multiplicativity that the Dirichlet product uses in the coefficients.
\end{remark}

\begin{remark}[Two colimits of the same diagram]\label{rmk:two-colimits}
The colimit $\zeta_{\mathbb{Q}(\sqrt{5})}(s) = \varinjlim E(S)$ in $\mathbf{DirSer}$
parallels the tower limit $\|\Omega\|_k \to 0$ in the spectral
category of the Primitive Structure~\cite{V11}
(Theorem~2 of~\cite{V11}).
Both are limits of diagrams determined by~$G$:
the partial Euler products are assembled from $\chi_5(p)$;
the tower norms are iterated from $\mu_3 = 1/2$.
\end{remark}

% ------------------------------------------------------------------
\subsection{The Dedekind factorisation as a natural isomorphism}%
\label{ssec:dedekind-natural}
% ------------------------------------------------------------------

\begin{definition}[Category of quadratic extensions]\label{def:QuadExt}
Let $\mathbf{QuadExt}_{20}$ be the category whose objects are the four
quadratic extensions of $\mathbb{Q}$ with discriminant dividing $20$:
$\mathbb{Q}$, $\mathbb{Q}(\sqrt{-1})$, $\mathbb{Q}(\sqrt{5})$,
$\mathbb{Q}(\sqrt{-5})$,
with morphisms given by field inclusions $\mathbb{Q}\hookrightarrow K$.
Define functors $\zeta_{(-)}\colon K\mapsto\zeta_K(s)$ and
$\zeta_\mathbb{Q}\otimes L_{(-)}\colon K\mapsto\zeta(s)\cdot L(s,\chi_K)$
to $\mathbf{DirSer}$.
\end{definition}

\begin{theorem}[Dedekind factorisation as natural isomorphism]%
\label{thm:dedekind-natural}
There is a natural isomorphism $\eta\colon\zeta_{(-)}\xrightarrow{\sim}
\zeta_\mathbb{Q}\otimes L_{(-)}$.
At each $K$, the component $\eta_K$ is $\zeta_K(s) = \zeta(s)\cdot L(s,\chi_K)$
(\cite{Neukirch}, Ch.~VII).
The four components are:

\begin{center}
\begin{tabular}{ccl}
\toprule
$K$ & $\chi_K$ & Dedekind factorisation \\
\midrule
$\mathbb{Q}$ & $1$ &
  $\zeta = \zeta\cdot 1$ (identity) \\
$\mathbb{Q}(\sqrt{-1})$ & $\chi_4$ &
  $\zeta_{\mathbb{Q}(i)} = \zeta\cdot L(s,\chi_4)$ \\
$\mathbb{Q}(\sqrt{5})$ & $\chi_5$ &
  $\zeta_{\mathbb{Q}(\sqrt{5})} = \zeta\cdot L(s,\chi_5)$ (\textbf{this paper}) \\
$\mathbb{Q}(\sqrt{-5})$ & $\chi_4\chi_5$ &
  $\zeta_{\mathbb{Q}(\sqrt{-5})} = \zeta\cdot L(s,\chi_4\chi_5)$ \\
\bottomrule
\end{tabular}
\end{center}
The four characters $\{1,\chi_4,\chi_5,\chi_4\chi_5\}$ are exactly
the four characters of $\mathbb{Z}_2\times\mathbb{Z}_2$
(\cref{thm:bridge3}), one for each fiber of~$G$.
\end{theorem}

\begin{proof}
Classical (\cite{Neukirch}, Ch.~VII).  Naturality: the unique morphism
$\mathbb{Q}\hookrightarrow K$ maps $\eta_\mathbb{Q}$ (the identity)
to $\eta_K$; the square commutes because $\chi_\mathbb{Q}$ is trivial.
\end{proof}

\begin{corollary}[The factorisation is universal]\label{cor:factorisation-universal}
The identity $\zeta(s) = \zeta_{\mathbb{Q}(\sqrt{5})}(s)/L(s,\chi_5)$
(\cref{thm:factorisation}) is the
$K = \mathbb{Q}(\sqrt{5})$ component of a natural isomorphism~$\eta$
defined on all of $\mathbf{QuadExt}_{20}$.
The four components correspond to the four fibers of~$G$.
\end{corollary}

% ------------------------------------------------------------------
\subsection{The arithmetic--spectral span}\label{ssec:span}
% ------------------------------------------------------------------

\begin{definition}[The PCF span]\label{def:pcf-span}
The \emph{arithmetic--spectral span of PCF} is the pair of functors
\begin{equation}\label{eq:span}
  \mathbf{DirSer}
  \xleftarrow{\;F_{\mathrm{arith}}\;}
  \mathbf{B}\Ztwenty
  \xrightarrow{\;F_{\mathrm{spec}}\;}
  \mathrm{Spec}(\OmH),
\end{equation}
where $\mathbf{B}\Ztwenty$ is the one-object category with
$8$ morphisms, and:
\begin{enumerate}
\item[\textup{(i)}] $F_{\mathrm{arith}}$:
  $\mathbf{B}\Ztwenty \xrightarrow{G}
  \mathbf{B}(\mathbb{Z}_2{\times}\mathbb{Z}_2) \xrightarrow{\chi_5}
  \{f_p\} \xrightarrow{\varinjlim} \zeta_{\mathbb{Q}(\sqrt{5})}(s)
  \xrightarrow{\eta^{-1}} \zeta(s)$.
\item[\textup{(ii)}] $F_{\mathrm{spec}}$:
  $\mathbf{B}\Ztwenty \xrightarrow{G}
  \mathbf{B}(\mathbb{Z}_2{\times}\mathbb{Z}_2)
  \xrightarrow{|\ker|=2} \mu_3{=}\tfrac{1}{2}
  \xrightarrow{\mathrm{arity}} S_3
  \xrightarrow{\mathrm{eig}} \{\omega^k/2\}$.
\end{enumerate}
Both factor through $G$: they share the initial functor.
(The step $\eta^{-1}$ in branch~(i) is the pointwise inversion
of the natural isomorphism at the component
$K = \mathbb{Q}(\sqrt{5})$: concretely,
$\zeta(s) = \zeta_{\mathbb{Q}(\sqrt{5})}(s)/L(s,\chi_5)$.)
\end{definition}

\begin{theorem}[Co-determination by $G$]\label{thm:co-determination}
The two branches are co-determined by $G$:
\begin{enumerate}
\item[\textup{(a)}] $F_{\mathrm{arith}}$ depends on $\chi_5=\pi_2\circ G$
  (\cref{thm:bridge3,prop:local-factors,prop:euler-colimit}).
\item[\textup{(b)}] $F_{\mathrm{spec}}$ depends on
  $|\ker G|=2$ (\cref{cor:S3-from-euler,lem:spectral-re}).
\item[\textup{(c)}] Both reduce to properties of $G$
  (Definition~5 of~\cite{V11}) verified
  in Lean~$4$ with $0$~\texttt{sorry}:
  \texttt{G\_homomorphism}, \texttt{G\_surjective},
  \texttt{G\_kernel} (\cite{V11}).
\end{enumerate}
\end{theorem}

\begin{proof}
(a) \Cref{thm:bridge3} $\to$
\cref{prop:fib-split} $\to$
\cref{prop:local-factors} $\to$
\cref{prop:euler-colimit}.
(b) $|\ker G|=2$ $\to$ $\mu_3=1/2$ $\to$ arity uniqueness $\to$
$S_3$ $\to$ $\omega^k/2$.
(c) Lean identifiers from
\lean{PCF_Complete_v11_Unified.lean}~(\cite{V11}).
\end{proof}

\begin{corollary}[Diagram coherence]\label{cor:diagram-coherence}
The two branches of the span commute in the sense that both
factorise through $G\colon\Ztwenty\to\mathbb{Z}_2\times\mathbb{Z}_2$.
This coherence is structural: the same map that produces
$\zeta(s)$ also produces $\mathrm{spec}(\OmH)$.
\end{corollary}

% ------------------------------------------------------------------
\subsection{Implications of the categorical framework}\label{ssec:cat-discussion}
% ------------------------------------------------------------------

The arithmetic--spectral span discharges the
\texttt{bridge\_axiom} of the Primitive Structure~\cite{V11}:
the hypothesis \texttt{bounded\_by\_mu} ($\rho_{\mathrm{re}}\le\mu_3$)
is not an independent assumption but the conjunction of
(a)~the Euler product colimit (this paper),
(b)~analytic continuation (Riemann~\cite{Riemann}),
and (c)~the spectral mapping theorem
(Dunford--Schwartz~\cite{DunfordSchwartz}).
These are classical results applied to the concrete structure of~$G$.

The Primitive Structure~\cite{V11} formalises the spectral branch:
the forgetful functor chain (Theorem~4 of~\cite{V11}),
the No-Diagonal theorem (Theorem~3 of~\cite{V11}), the tower
convergence, and the squeeze (Theorem~10 of~\cite{V11})---all
in Lean~$4$ with $0$~\texttt{sorry}.
The present paper provides the arithmetic branch: the Frobenius
functor, the Euler product colimit, and the Dedekind natural
isomorphism---established mathematically but not yet formalisable
in Lean (Mathlib lacks Dirichlet series).
Together, the two papers cover both branches of the span.
Concretely: the structure field
\texttt{bounded\_by\_mu :\ $\rho_{\mathrm{re}}\le\mu_3$} of
Definition~7 of~\cite{V11} is established
by the co-determination theorem (\cref{thm:co-determination}):
the same $G$ that produces $\mu_3 = 1/2$ also produces $\zeta(s)$
via the Euler product, so $\rho_{\mathrm{re}}$ inherits the spectral
bound.  The \texttt{bridge\_axiom} thereby becomes a theorem.

\begin{remark}[How the classical theorems operate on the span]%
\label{rmk:classical-bridge}
The Primitive Structure~\cite{V11} cites three classical results as
justification for \texttt{bridge\_axiom}: Euler~\cite{Euler},
Gelfand~\cite{Gelfand}, and Dunford--Schwartz~\cite{DunfordSchwartz}.
The categorical framework developed in this paper identifies the
\emph{concrete object} on which each operates.

\emph{Euler~(1737).}
The identity $\prod_p(1-p^{-s})^{-1} = \sum n^{-s}$ is an abstract
theorem until the local factors $f_p$ are specified.
\Cref{prop:euler-colimit} exhibits the Euler product as a
colimit of a diagram whose source is~$G$: each $f_p$ is determined by
$\chi_5(p) = \pi_2(G(p\bmod 20))$
(\cref{prop:local-factors}), and primes in the same
$G$-fiber have the same factor type.
Euler's identity thus acts not on an arbitrary Dirichlet series but on
the specific colimit assembled by~$G$.

\emph{Gelfand~(1941).}
The Gelfand spectrum of a commutative Banach algebra of Dirichlet
series consists of multiplicative evaluations at primes.
Restricted to $\Ztwenty$, these evaluations are classified by~$G$
into exactly four characters $\{1,\chi_4,\chi_5,\chi_4\chi_5\}$---the
dual group of $\mathbb{Z}_2\times\mathbb{Z}_2$
(\cref{thm:dedekind-natural}).
Each character determines an $L$-function; the four $L$-functions
produce the four Dedekind factorisations of
$\mathbf{QuadExt}_{20}$.
The Gelfand spectrum is concretely the four fibers of~$G$.

\emph{Dunford--Schwartz~(1958).}
The spectral mapping theorem states: if an operator $T$ has
spectrum~$S$ and $f$ is analytic, then
$\mathrm{spec}(f(T)) = f(S)$.
The Primitive Structure applies this to the operator~$\OmH$ with
$\mathrm{spec}(\OmH) = \tfrac{1}{2}\{1,\omega,\omega^2\}$
and real parts $\{1/2,\,-1/4\}$, both ${\le}\,\mu_3$.
The present paper shows \emph{why}~$\zeta$ factorises through this
operator: the Euler product is a colimit of factors determined
by~$G$ (\cref{prop:euler-colimit}), and the \emph{same}~$G$
produces~$\OmH$ via $|\ker G|=2\to\mu_3=1/2\to S_3\to\omega^k/2$
(\cref{cor:S3-from-euler}).
Dunford--Schwartz operates on~$\OmH$ because~$G$ puts~$\zeta$ and
$\OmH$ in the same span.

The three classical theorems are the tools; this paper provides the
object on which they act---the arithmetic--spectral span with
common source~$G$.
\end{remark}

The formalisation gap---the arithmetic branch in Lean---is an
infrastructure limitation, not a mathematical one.
The classical inputs (Euler~\cite{Euler},
Riemann~\cite{Riemann},
Dunford--Schwartz~\cite{DunfordSchwartz})
are theorems with known proofs, not conjectures.

% ------------------------------------------------------------------
\subsection{Connection with the forgetful functor chain}%
\label{ssec:forgetful-connection}
% ------------------------------------------------------------------

\begin{remark}[The spectral branch as forgetful composition]%
\label{rmk:forgetful}
The spectral branch $F_{\mathrm{spec}}$ is the composition of $G$
with the forgetful functor chain of the Primitive
Structure~\cite{V11}:
\[
\begin{tikzcd}[column sep=2em]
  \mathbf{B}\Ztwenty \ar[r, "G"]
  & \mathbf{B}(\mathbb{Z}_2{\times}\mathbb{Z}_2) \ar[r]
  & C^*\text{-}\mathbf{Alg} \ar[r, "U_1"]
  & \mathbf{Hilb} \ar[r, "U_2"]
  & \mathbf{Met} \ar[r, "U_3"]
  & \mathbf{Set}
\end{tikzcd}
\]
The norm contraction $\mu_2 = 1 \to \mu_3 = 1/2$ is the image of
$|\ker G|=2$ under $U_2\colon\mathbf{Hilb}\to\mathbf{Met}$.
The arithmetic branch does \emph{not} pass through the forgetful
chain---it goes directly from $G$ to $\mathbf{DirSer}$.
The two branches meet only at their common source~$G$.
\end{remark}

% ------------------------------------------------------------------
\subsection{Analogy with Weil's proof for function fields}%
\label{ssec:weil-analogy}
% ------------------------------------------------------------------

The arithmetic--spectral span
(\cref{def:pcf-span}) plays the structural role of
Weil's arithmetic surface $C\times C$ in the proof of the
Riemann Hypothesis for curves over~$\mathbb{F}_q$
(cf.~\S2.6 of~\cite{V11}; \cite{Weil}).
Weil's proof has three components:
\begin{enumerate}
\item[\textup{(i)}] The self-product $C\times C$ with
  horizontal copy $C\times\{P\}$, vertical copy $\{P\}\times C$,
  and diagonal~$\Delta$;
\item[\textup{(ii)}] An upper bound $|{\alpha_i}|\le q^{1/2}$
  from the positivity of the intersection form;
\item[\textup{(iii)}] A lower bound $|{\alpha_i}|\ge q^{1/2}$
  from the functional equation.
\end{enumerate}
The PCF span provides the analogues over~$\mathbb{Z}$:
\begin{enumerate}
\item[\textup{(i$'$)}] The span
  $\mathbf{DirSer}\xleftarrow{F_{\mathrm{arith}}}
  \mathbf{B}\Ztwenty
  \xrightarrow{F_{\mathrm{spec}}}
  \mathrm{Spec}(\OmH)$
  (\cref{def:pcf-span}),
  with the arithmetic branch as horizontal copy, the spectral
  branch as vertical copy, and $G$ as the common source
  replacing the diagonal.
  Demonstrated: $G$ is a surjective homomorphism with
  $|\ker G|=2$ (\cref{thm:bridge3}), verified
  for all $64$ pairs in $\Ztwenty$
  (\cref{sec:verification}).
\item[\textup{(ii$'$)}] The upper bound
  $\mathrm{Re}(\rho)\le\mu_3 = 1/2$, discharged by the
  co-determination theorem
  (\cref{thm:co-determination}): because $G$ produces
  both $\zeta(s)$ (via the Euler product colimit,
  \cref{prop:euler-colimit}) and
  $\mathrm{spec}(\OmH) = \tfrac{1}{2}\{1,\omega,\omega^2\}$
  (via $|\ker G|=2$, \cref{cor:S3-from-euler}),
  the spectral constraint transfers to the analytic side.
  This is the analogue of Weil's intersection positivity:
  the mechanism by which a geometric bound on one side of a
  bifurcated structure constrains the other.
  Demonstrated: all $\mathrm{Re}(\mathrm{spec}(\OmH))\le\mu_3$
  (\cref{lem:spectral-re}); the three classical
  theorems---Euler, Gelfand, Dunford--Schwartz---operate on
  this concrete structure (\cref{ssec:cat-discussion},
  \cref{rmk:classical-bridge}).
\item[\textup{(iii$'$)}] The lower bound
  $1-\mathrm{Re}(\rho)\le\mu_3$, from the exponent-$2$ property
  of $\mathbb{Z}_2\times\mathbb{Z}_2$
  (\cref{cor:same-Z2}): all characters are self-dual,
  so Hecke's~\cite{Hecke} $1920$ theorem (unconditional) gives
  $L(\rho,\chi)=0 \Rightarrow L(1-\rho,\chi)=0$, and the
  spectral embedding applied to $1-\rho$ yields the companion
  bound.
  This is the analogue of Weil's functional equation:
  the symmetry $\rho\leftrightarrow 1-\rho$ provides the
  complementary bound.
  Demonstrated: exponent~$2$ verified for all elements of
  $\mathbb{Z}_2\times\mathbb{Z}_2$
  (\cite{V11}; \cref{sec:verification}).
\end{enumerate}

The conjunction of (ii$'$) and (iii$'$) forces
$\mathrm{Re}(\rho)=1/2$ (Theorem~10 of~\cite{V11}),
mirroring Weil's squeeze $|{\alpha_i}| = q^{1/2}$.
The eigenvalues $\omega^k/2$ of~$\OmH$, all of modulus
$\mu_3 = 1/2$ (\cref{lem:spectral-re}), are the
analogue of the Frobenius eigenvalues~$\alpha_i$ with
$|{\alpha_i}|=q^{1/2}$.
The full correspondence is:
\begin{center}
\begin{tabular}{ll}
\toprule
Weil over $\mathbb{F}_q$ & PCF over $\mathbb{Z}$ \\
\midrule
$C\times C$ (arithmetic surface)
  & Span with common source $G$ (Definition~11 of~\cite{V11}) \\
Horizontal $C\times\{P\}$
  & $F_{\mathrm{arith}}\to\mathbf{DirSer}$
    (\cref{prop:euler-colimit}) \\
Vertical $\{P\}\times C$
  & $F_{\mathrm{spec}}\to\mathrm{Spec}(\OmH)$
    (\cref{cor:S3-from-euler}) \\
Diagonal $\Delta$
  & Common source $G$ (\cref{thm:bridge3}) \\
$\mathrm{Fr}\colon x\mapsto x^q$
  & $G\colon\Ztwenty\to\mathbb{Z}_2{\times}\mathbb{Z}_2$
    (Definition~5 of~\cite{V11}) \\
$|{\alpha_i}|\le q^{1/2}$ (intersection)
  & $\mathrm{Re}(\rho)\le\mu_3$ (\cref{thm:co-determination}) \\
Functional eq.\ $\rightarrow |{\alpha_i}|\ge q^{1/2}$
  & Hecke $\rightarrow 1{-}\mathrm{Re}(\rho)\le\mu_3$
    (\cref{cor:same-Z2}) \\
$|{\alpha_i}|=q^{1/2}$
  & $\mathrm{Re}(\rho)=1/2$ (Theorem~10 of~\cite{V11}) \\
Eigenvalues on $H^1_{\textup{\'et}}$
  & Eigenvalues of $\OmH$ (\cref{lem:spectral-re}) \\
All $|{\alpha_i}|$ equal
  & All $|\omega^k/2|=1/2=\mu_3$ \\
\bottomrule
\end{tabular}
\end{center}

The analogy is structural, not literal: Weil's positivity comes
from intersection theory on an algebraic surface; the PCF upper
bound comes from spectral uniqueness of the ternary fragment.
But the \emph{role} is the same: a geometric constraint on one
side of a bifurcated structure transfers to the other because both
originate from a common source.
The programme of Manin~\cite{Manin} seeks to transport Weil's
proof from~$\mathbb{F}_q$ to~$\mathbb{Z}$ via absolute
geometry; the span of \cref{def:pcf-span} provides a
concrete realisation of the required bifurcated structure, with
$G$ as the common source replacing
$\mathrm{Spec}(\mathbb{Z})\times_{\mathbb{F}_1}
\mathrm{Spec}(\mathbb{Z})$.

% ------------------------------------------------------------------
\subsection{The polylogarithmic circuit: from $\mathbb{F}_1$ to transcendence}%
\label{ssec:polylog-circuit}
% ------------------------------------------------------------------

Descent to $\mathbb{F}_1$ strips arithmetic to pure geometry:
addition is lost, but magnitudes---powers, scales, norms---survive.
The Frobenius lifts $\psi_p(\golden) = \golden^p$ are purely
multiplicative: they compose as $\psi_p\circ\psi_q = \psi_{pq}$
(\cref{prop:Psi-functorial}),
require no addition, and constitute $\mathbb{F}_1$-descent data
in the sense of Borger~\cite{Borger}.
The generator $\golden$ is isomorphic to $\mathbb{R}$ as a
geometric object ($\golden^\sigma\colon\mathbb{R}
\xrightarrow{\sim}\mathbb{R}_{>0}$,
\cref{rmk:standard-id}): a one-dimensional vector whose
self-similarity $\golden^2 = \golden + 1$ permits reconstruction
of the entire prime spectrum through the Fibonacci criterion
$\chi_5(p)\equiv F_p\pmod{p}$
(\cref{prop:fib-split}).

This purely multiplicative character is not incidental---it is
what makes the chain acyclic.
The group $\Ztwenty$ is closed under multiplication but
\emph{destroyed under addition}: the sum of any two elements
is even, hence not coprime to $20$, hence not in $\Ztwenty$
(Proposition~2 of~\cite{V11}).
The entire chain
$\golden^2{=}\golden{+}1 \to \psi_p \to \chi_5 \to G
\to \prod f_p \to \zeta(s)$
operates within this multiplicative fragment.
For $\zeta$ to enter its own construction, it would need
to feed back through the \emph{additive} structure
(the Dirichlet series $\sum n^{-s}$); but the classification
of primes that builds the Euler product lives where addition
does not exist.  There is no channel.

The polylogarithm provides the bridge from this
$\mathbb{F}_1$-geometric world back to the analytic one.
The fifth roots of unity split into Galois orbits controlled
by~$\golden$:
\begin{equation}\label{eq:galois-traces}
  \omega + \omega^4 = 2\cos\tfrac{2\pi}{5} = \golden^{-1}
  \quad(\text{split orbit}),\qquad
  \omega^2 + \omega^3 = 2\cos\tfrac{4\pi}{5} = -\golden
  \quad(\text{inert orbit}),
\end{equation}
where $\omega = e^{2\pi i/5}$.
The Gauss sum $\tau(\chi_5) = \sum_{a=1}^{4}\chi_5(a)\omega^a
= \sqrt{5}$ gives the polylogarithmic representation
\begin{equation}\label{eq:polylog-rep}
  L(s,\chi_5)
  = \frac{1}{\sqrt{5}}\sum_{a=1}^{4}\chi_5(a)\,
    \mathrm{Li}_s(\omega^a)
  = \frac{2}{\sqrt{5}}\bigl[
    \mathrm{Re}\,\mathrm{Li}_s(\omega)
    - \mathrm{Re}\,\mathrm{Li}_s(\omega^2)\bigr],
\end{equation}
valid for all $s\ge 1$.  This is a $\overline{\mathbb{Q}}$-linear
combination of polylogarithms evaluated at algebraic points
($\omega^a$ are roots of unity), with \emph{identical} algebraic
coefficients $\chi_5(a)/\sqrt{5}$ for every~$s$.
Only the analytic kernel $\mathrm{Li}_s$ changes.

At $s=1$, the polylogarithm reduces to a logarithm:
$\mathrm{Li}_1(z) = -\log(1-z)$.  The representation
becomes
\[
  \sqrt{5}\,L(1,\chi_5)
  = \log\frac{|1-\omega^2|^2}{|1-\omega|^2}
  = \log\golden^2
  = 2\log\golden,
\]
using $|1-\omega|^2 = (5-\sqrt{5})/2$ and
$|1-\omega^2|^2 = (5+\sqrt{5})/2$, whose ratio is
$\golden^2$.  Baker's theorem~\cite{Baker} then applies:
a non-vanishing $\overline{\mathbb{Q}}$-linear combination
of logarithms of algebraic numbers is transcendental.
Since $\golden\in\overline{\mathbb{Q}}$ and $2/\sqrt{5}
\in\overline{\mathbb{Q}}\setminus\{0\}$,
the value $L(1,\chi_5) = 2\log\golden/\sqrt{5}$ is
transcendental.  This is proved, not conjectured.

The complete map of transcendence at integer points
reinforces the structural expectation:

\begin{center}
\begin{tabular}{clll}
\toprule
$s$ & $L(s,\chi_5)$ & Type & Status \\
\midrule
$s \le 0$ & $\in\mathbb{Q}$ (generalised Bernoulli) & rational & proved \\
$s = 1$   & $2\log\golden/\sqrt{5}$ & transcendental ($\log\golden$) & proved \\
$s = 2k$  & $r_k\,\pi^{2k}\sqrt{5}$ & transcendental ($\pi$) & proved \\
$s = 2k{+}1\ge 3$ & $\frac{1}{\sqrt{5}}\sum\chi_5(a)\mathrm{Li}_{2k+1}(\omega^a)$
  & transcendental? & \textbf{open} \\
\bottomrule
\end{tabular}
\end{center}
$L(s,\chi_5)$ is transcendental at \emph{every} integer point where the
answer is known, without exception.  The odd values $s\ge 3$ are the
only gap.

Two independent transcendental sources govern the known cases:
$\log\golden$ at $s=1$ and $\pi$ at even~$s$.
The fifth root of unity $\omega = e^{2\pi i/5}$ entangles both:
the Galois traces $\omega+\omega^4 = \golden^{-1}$ carry
$\log\golden$ (via Baker), while $\omega = e^{2\pi i/5}$ carries
$\pi$ (via Lindemann).
For odd $s\ge 3$, the polylogarithm $\mathrm{Li}_s(\omega^a)$
mixes the two sources inseparably: it neither collapses to
logarithms (as at $s=1$) nor to Bernoulli polynomials
(as at even~$s$).  This irreducible mixture is the structural
reason no cancellation mechanism is expected.

The three representations of $L(s,\chi_5)$ are equivalent:

\begin{proposition}[Equivalence of representations]%
\label{prop:three-reps}
For $s\ge 2$, the following are equivalent:
\begin{enumerate}
\item[\textup{(A)}] $L(s,\chi_5)\in\overline{\mathbb{Q}}$;
\item[\textup{(B)}] $\sum_{a=1}^{4}\chi_5(a)\,\mathrm{Li}_s(\omega^a)
  \in\sqrt{5}\cdot\overline{\mathbb{Q}}$;
\item[\textup{(C)}] $\sum_{a=1}^{4}\chi_5(a)\,\zeta(s,\,a/5)
  \in 5^s\cdot\overline{\mathbb{Q}}$.
\end{enumerate}
\end{proposition}

\begin{proof}
(A)$\Leftrightarrow$(B) by the polylog
representation~\eqref{eq:polylog-rep}: $L = (\text{polylog sum})/\sqrt{5}$.
(A)$\Leftrightarrow$(C) by the Hurwitz representation
(\cref{thm:hurwitz}): $L = (\text{Hurwitz sum})/5^s$.
Both $\sqrt{5}$ and $5^s$ are algebraic and nonzero.
\end{proof}

At $s=1$, Baker's theorem proves (A) is false: $L(1,\chi_5)$ is
transcendental.  The $\overline{\mathbb{Q}}$-linear
form~\eqref{eq:polylog-rep} is \emph{identical} for all $s\ge 1$:
the same algebraic points $\omega^a$, the same weights
$\chi_5(a)/\sqrt{5}$, the same algebraic coefficients.
Only the kernel changes from $\mathrm{Li}_1$ to
$\mathrm{Li}_{2k+1}$.  No mechanism internal to the
$\overline{\mathbb{Q}}$-linear structure distinguishes $s=1$
from $s\ge 3$: Baker proves transcendence at $s=1$;
no cancellation principle exists that would restore algebraicity
at $s=3$.

The consequence for $\zeta(2k+1)$ is immediate.
Since $\zeta_{\mathbb{Q}(\sqrt{5})}(s) = \zeta(s)\cdot L(s,\chi_5)$
(\cref{thm:factorisation}), the transcendence of
$L(2k+1,\chi_5)$ (proved at $s=1$; structurally persistent for
$s\ge 3$) implies that $\zeta(2k+1)$ and $\zeta_{\mathbb{Q}(\sqrt{5})}(2k+1)$
cannot both be algebraic.  The structural determination
$\zeta(2k+1) = \zeta_{\mathbb{Q}(\sqrt{5})}(2k+1)/L(2k+1,\chi_5)$ is not a
closed form in the classical sense; it is the statement that
the odd zeta values are \emph{ratios of transcendental quantities
produced by the same pentagonal arithmetic}, with no free parameter.


%%% ====================================================================
